%% chapters/14-goe.tex
\chapter{The Gaussian Orthogonal Ensemble (GOE)}
\label{ch:goe}

\whereWeAre{%
The non-Gaussian limit law from Chapter 12 has a name: the eigenvalue distribution of the \emph{Gaussian Orthogonal Ensemble}. This chapter formalizes the GOE, derives its joint eigenvalue density, and explains the Vandermonde repulsion factor that makes coincident eigenvalues vanishingly unlikely.%
}

\PLACEHOLDER{Sources: existing main.tex GOE section; Joseph thesis Section 3.2.}

\section{Definition of the GOE}
\PLACEHOLDER{$N \times N$ symmetric, diagonal entries $\Normal(0,2)$, off-diagonal $\Normal(0,1)$, all independent up to symmetry. Construction $H = (G+G^T)/\sqrt{2}$ from i.i.d.\ Gaussian $G$.}

\section{Joint eigenvalue density and the Vandermonde}
\PLACEHOLDER{Density formula with $\prod_{i<j} |\lambda_i - \lambda_j|^{\beta=1}$ factor; explicit $2\times 2$ case.}

\section{Eigenvalue repulsion}
\PLACEHOLDER{The Vandermonde forces the density to vanish on the diagonal $\lambda_1 = \lambda_2$; eigenvalues ``avoid'' each other.}

\section{Why GOE specifically appears in our story}
\PLACEHOLDER{$\sqrt{n}(\Shat - I_p)$ converges to a GOE-like matrix in fixed dimension. Closes the loop with Chapter 12.}

\section{Brief warm-ups}
\PLACEHOLDER{3--5 short questions.}

\chapterRecap{%
\begin{itemize}
  \item \PLACEHOLDER{GOE = $\beta = 1$ symmetric Gaussian ensemble.}
  \item \PLACEHOLDER{Vandermonde factor encodes repulsion.}
  \item \PLACEHOLDER{Limit law for $\sqrt n(\Shat - I_p)$ in fixed $p$ is GOE-shaped.}
\end{itemize}%
}

\section*{Exercises}
\PLACEHOLDER{8--15 longer exercises.}
